\chapter{A Venue Playbook}\label{ch:hf:a-venue-playbook}

The same quoting code behaves differently on a price-time equities exchange, a pro-rata futures book, a venue with a \omterm{def:hf:latency-arbitrage-and-its-defence:bump}{speed bump} and a crypto
exchange with an order budget. A desk that learns the differences in production learns them from its losses: orders rejected by a budget it did not
read, a queue it joined with the wrong size, quotes that arrive five seconds after the price it saw. So a desk keeps one page per venue, and checks
the page before the first order. In this chapter's simulation, one quoter written once and run with its home venue's parameters has 94\% of its
orders rejected on a crypto venue, trades a quarter of the volume it could on a pro-rata book, and turns half of its volume into takes on a venue that
delays new orders. Three venues need a parameter changed, four need nothing, and two, one of them among the three, need different code.

\section{What a venue profile holds}

\begin{definition}[Venue profile]\label{def:hf:a-venue-playbook:profile}
A \emph{venue profile}\index{venue profile} is a structured, versioned record of what a trading strategy must know about a venue before sending
it an order (its matching rule, tick and lot, fees and their unit, market-data feed, order types, delays, message budget, session and known quirks),
together with the strategy's own parameters for that venue, validated by rules that tie the two together.
\end{definition}

A profile has two halves (\cref{fig:hf:a-venue-playbook:flow}). The first describes the venue and changes when the venue changes: a new fee
schedule, a new order type, a new limit. The second describes how the desk quotes there and changes when the desk learns: the size it shows, how far
the price must move before it requotes, how often it may requote, whether it flags its orders post-only. The rules tie the halves together. A
quote size must be a multiple of the lot. A post-only flag needs a venue that offers post-only orders. A delay that spares post-only orders
should be met with post-only orders. Requotes must fit the message budget. Each rule is one line of \texttt{firm.venueprofile.validate}
(\cref{lst:hf:a-venue-playbook:validate}), and each was learned by somebody in production.

The fields, one by one, are what the earlier chapters used:
\begin{itemize}
\item \textbf{Matching.} Price-time priority (chapters 1 to 10) or pro-rata allocation with a top-order share (\cref{ch:hf:futures-market-making}),
or a batch interval with one clearing price (\cref{ch:hf:latency-arbitrage-and-its-defence}), or no matching engine at all
(\cref{ch:hf:on-chain-trading}).
\item \textbf{Tick, lot and price scale.} A tick of one cent on a \$100 stock is one basis point; a quarter point on a futures contract is a
different amount of money per contract, and the lot turns a quote size into risk.
\item \textbf{Fees and their unit.} Per share on US equities, per contract on futures and options, in basis points of notional on crypto venues;
the sign of the maker fee separates a maker-taker venue from an inverted one (\cref{ch:hf:rebates-inverted-venues-and-tiers}).
\item \textbf{Feed.} Market-by-order or market-by-price, and whether a quoter can find its own orders in it: a quoter that cannot will chase them
(\cref{ch:hf:the-business-of-market-making}).
\item \textbf{Order types.} Post-only, \omterm{def:hf:automated-options-market-making:mass}{mass quotes}, replace with or without loss of priority, pegs.
\item \textbf{Delays.} A \omterm{def:hf:latency-arbitrage-and-its-defence:bump}{speed bump}, symmetric or asymmetric (\cref{def:hf:latency-arbitrage-and-its-defence:bump}); an in-play delay on
placing orders; a batch interval.
\item \textbf{Budgets.} Orders a second, a day, per account and per instrument, whether cancels count, and what happens on a breach: a reject, a
disconnection, a fee.
\item \textbf{Session and quirks.} Hours, auctions, maintenance windows, and the paragraph that begins ``note that''.
\end{itemize}

\begin{omfigure}
\begin{tikzpicture}[font=\footnotesize,>=stealth,box/.style={draw,rounded corners,fill=omThm!8,minimum width=2.5cm,minimum height=0.75cm,align=center}]
\node[box] (v) at (0,1.3) {venue half\\matching, tick, lot, fees,\\feed, orders, delays, budget};
\node[box] (q) at (0,-0.3) {quoting half\\size, requote threshold,\\pacing, post-only};
\node[box,fill=omDef!12] (val) at (4.2,0.5) {validate\\(rules tying\\the halves)};
\node[box] (cfg) at (8.2,1.3) {simulator\\configuration};
\node[box] (par) at (8.2,-0.3) {quoter\\parameters};
\node[box,fill=omDef!20] (sim) at (8.2,-1.9) {one quoter,\\home vs own};
\node[box,fill=gray!10] (rep) at (4.2,-1.9) {report: volume, edge,\\picked off, messages};
\node[box,fill=gray!10] (reg) at (0,-1.9) {registry row\\and fee schedule};
\draw[->] (v) -- (val); \draw[->] (q) -- (val);
\draw[->] (val) -- (cfg); \draw[->] (val) -- (par);
\draw[->] (cfg.east) -- ++(0.5,0) |- (sim.east); \draw[->] (par) -- (sim);
\draw[->] (sim) -- (rep); \draw[->] (val.south west) -- (reg.north east);
\end{tikzpicture}

\omcaption{A \omterm{def:hf:a-venue-playbook:profile}{venue profile} and what it produces. The venue half and the quoting half are validated together; a valid profile yields the
simulator's configuration (Book 10's exchange simulator), the quoter's parameters, and the rows for Book 1's venue registry and fee-schedule
engine. The chapter runs one quoter on every simulated profile with the home venue's parameters and with the profile's own. Source:
\texttt{firm.venueprofile}.}\label{fig:hf:a-venue-playbook:flow}
\end{omfigure}

The chapter's registry holds eleven profiles, one per venue type; eight have a matching engine the exchange simulator can run. Their fees and
limits are illustrative of each type, set from the dated box's figures where one exists; a desk's own profile holds the venue's actual schedule, with
the date it was read. The venue half of the eleven:

\begin{center}
\footnotesize
\setlength{\tabcolsep}{4pt}
\begin{tabular}{@{}llllll@{}}
\toprule
profile & matching & tick; lot & fees, make/take & delay or budget & own quoting\\
\midrule
lit equity & price-time & 1 c; 100 & $-0.20$/$+0.30$ c/sh & -- & home\\
inverted equity & price-time & 1 c; 100 & $+0.10$/$-0.05$ c/sh & -- & home\\
bumped equity & price-time & 1 c; 100 & $-0.20$/$+0.30$ c/sh & 350 \textmu s, all & home\\
batch & batch auction & 1 c; 100 & 0/0 & 100 ms interval & home\\
pro-rata futures & top, pro rata & 0.25; 1 & 0/0 a contract & -- & size 500\\
FIFO futures & price-time & 0.25; 1 & 0/0 a contract & message ratio & home\\
crypto & price-time & 1; 1 & $-0.5$/$+2.5$ bp & 0.125/s, burst 5 & 16-s pacing\\
event & price-time & 1 c; 1 & 0/0 & 5 s, new orders & 3-tick threshold\\
options & top, pro rata & 1 c; 1 & $-25$/$+50$ c/ct & risk tools & not run\\
FX ECN & price-time & 1 pip; 1M & 0/0.05 bp & -- & not run\\
on-chain & pool & -- & 30 bp pool fee & block time & not run\\
\bottomrule
\end{tabular}
\end{center}

The crypto budget is a per-instrument share of an account's: 200\,000 orders a day, the dated box's figure, shared by about eighteen instruments,
leaves one order every eight seconds for each, with a burst of five. The event venue's five-second delay on new orders sits inside the one-to-twelve
seconds the dated box reports for one betting exchange's in-play markets.

\section{Equities and options venues}

On a US equity exchange the profile's first question is the fee sign. A maker-taker venue pays the maker a rebate and charges the taker up to the
access fee cap; an inverted venue charges the maker and pays the taker, so its queues are shorter and its fills are more often the informed ones
(\cref{ch:hf:rebates-inverted-venues-and-tiers}). In the chapter's run the two venues are the same book with the fee signs swapped: the quoter trades
14\,567 units on each, the same fills, and its fee line moves from a rebate of \$29 to a charge of \$15 over the twenty minutes. Nothing in the quoter's
parameters has to change for that; its accounting does, and so does the decision of whether to quote there at all.

The second question is the feed. An order-by-order feed (a market-by-order feed such as Nasdaq's ITCH, where a replaced order takes a new reference)
lets the quoter find its own orders and quote off everyone else's; a market-by-price feed shows totals per level, and a quoter must subtract its own
size, which is exact only once its acknowledgements have arrived. The chapter's quoter identifies its orders by the references in their
acknowledgements, and the simulator delivers acknowledgements before the public feed (15 against 20 microseconds), as private order entry
usually does. With the acknowledgement arriving second, a first draft of this quoter took its own order for somebody else's and chased it.

The third question is the delay. The bumped equity venue delays every message by 350 microseconds, the delay one US exchange has run since 2016.
In the run it changes nothing: the arbitrageur's order and the quoter's cancel are delayed alike, the arbitrageur still wins every race, and 42\% of
the quoter's volume is still taken on a jump. The bump in that design protects the venue's pegged orders, which reprice before the delayed
arbitrageur arrives (\cref{ch:hf:latency-arbitrage-and-its-defence}); it does not protect a displayed quote whose owner reacts through the same
delay. An asymmetric delay, one that spares cancels, would; a US proposal of that kind was disapproved in 2020.

\begin{omfigure}
\begin{tikzpicture}
\begin{axis}[width=11cm,height=5.4cm,ybar,bar width=9pt,ymin=0,ymax=72,xmin=-0.6,xmax=7.6,
  xtick={0,1,2,3,4,5,6,7},xticklabels={lit,inverted,bumped,batch,{pro rata},FIFO fut.,crypto,event},
  xticklabel style={font=\scriptsize,rotate=25,anchor=north east},ylabel={\small volume picked off (\%)},
  tick label style={font=\footnotesize},grid=major,grid style={gray!25},table/col sep=comma,
  legend style={font=\scriptsize,at={(0.03,0.97)},anchor=north west},legend columns=2]
\addplot[fill=omThm!60,draw=omThm] table[x=i,y=picked_home]{figdata/hft/29-a-venue-playbook/playbook.csv};
\addplot[fill=omDef!55,draw=omDef] table[x=i,y=picked_own]{figdata/hft/29-a-venue-playbook/playbook.csv};
\legend{home parameters,own parameters}
\end{axis}
\end{tikzpicture}

\omcaption{The share of the quoter's volume taken by the fast arbitrageur within a jump of the efficient price, on each simulated venue, with the
home venue's parameters and with the venue's own; three seeds of 20 minutes, a jump every 20 seconds on average. The arbitrageur hears of a jump 20
microseconds before the quoter and sends through a 5-microsecond link against the quoter's 20. Data: \texttt{hf\_venues.playbook}.}\label{fig:hf:a-venue-playbook:picked}
\end{omfigure}

Options venues add two fields a stock venue does not have: \omterm{def:hf:automated-options-market-making:mass}{mass quotes}, one message that replaces the bid and offer of many series, and
exchange-side quote risk protection, which stops a maker's executions when those across a class in a short window exceed limits it has set
(\cref{ch:hf:automated-options-market-making}). The profile records both, with the protection's parameters as the desk has set them; they are the
last line of defence on a venue where a single sweep touches hundreds of series.

\section{Futures and FX venues}

A pro-rata futures book shares each incoming order among the resting orders in proportion to their sizes, after giving a share to the order that
set the price. Size is the parameter: the quoter with the home venue's 100 units traded 10\,884 units in the runs, and with the profile's 500 it traded
40\,633, 3.7 times as much with fewer messages. The larger size also has more of itself taken on each jump, 58\% of its volume against 47\%, so the
edge per unit barely moves, $-0.14$ ticks against $-0.17$; whether to show the size depends on the edge, and on how many of the other makers
over-quote in turn (\cref{ch:hf:futures-market-making}). A FIFO futures book is a price-time book with contract units: in the runs it trades
exactly as the lit equity venue, fee line aside. What its profile adds is the exchange's message-efficiency benchmark, messages over traded volume
per product group, which a requote threshold spends (\cref{ch:hf:the-quoting-engine}).

\begin{dated}{2026-09}{One venue per type, as their rules read}\label{dat:hf:a-venue-playbook:venues}
\small
\begin{tabular}{@{}p{2.3cm}p{9.9cm}@{}}
US equities & Nasdaq's ITCH feed is order by order; a replaced order gets a new reference. The access fee cap is \$0.003 a share until the first
business day of November 2027, when the \$0.001 cap adopted in 2024 applies.\\
\omterm{def:hf:latency-arbitrage-and-its-defence:bump}{Speed bump} & IEX's approval as an exchange, on 17 June 2016, came with a 350-microsecond delay on its inbound and outbound messages. The SEC
disapproved an asymmetric four-millisecond delay on Cboe EDGA in February 2020.\\
Futures & CME's pro-rata algorithms give the order that first improves the market its share, then allocate pro rata, rounding down, with the
remainder FIFO (a vendor's guide, August 2025). CME's messaging efficiency programme compares messages over volume with product-group
benchmarks.\\
Options & Cboe's risk tools stop a maker's executions once its notional, volume, count or percentage-of-quote limit over a period of 100
milliseconds to a day is reached.\\
FX & Some primary venues introduced \omterm{def:hf:latency-arbitrage-and-its-defence:bump}{speed bumps} to protect dealers (BIS, 2022); last look must be disclosed under the FX Global Code.\\
Crypto & A large spot venue's endpoint returned order limits of 100 per 10 seconds and 200\,000 a day on 24 September 2026.\\
Betting & One betting exchange delays placing bets in play by one to twelve seconds so that customers can cancel unmatched bets.
\end{tabular}
\end{dated}

FX venues split the profile in two: a central limit order book (an ECN) with a firm-price protocol, and the disclosed or undisclosed last look of
the streaming channels (\cref{ch:hf:fx-market-making}). The profile for an ECN holds the minimum size (a million of the base currency
is a common lot, so the home venue's 100 units is not even a valid order: \texttt{validate} says so), the fee in basis points or per million, the
venue's latency floor or \omterm{def:hf:latency-arbitrage-and-its-defence:bump}{speed bump} if it has one, and the session, which runs five days round the clock with a weekly gap.

\section{Crypto, on-chain and event venues}

A crypto venue meters orders. With the profile's budget, one order every eight seconds and a burst of five, the quoter with the home venue's
parameters sent 2\,735 orders in twenty minutes and had 2\,581 of them rejected: every rejected order was sent again at the next book update, and each
rejection consumed nothing but left the quote missing. Paced to one requote every sixteen seconds, the rule the profile's validation enforces (two
orders a requote within the budget), it sent 118 orders, none rejected, and traded 8\,333 units against 7\,267
(\cref{fig:hf:a-venue-playbook:budget}).

\begin{omfigure}
\begin{tikzpicture}
\begin{axis}[width=11cm,height=5.4cm,xlabel={\small minimum gap between requotes (s)},ylabel={\small orders in 20 minutes},ymode=log,log basis y=10,
  ymin=10,ymax=5000,xmin=-1,xmax=33,xtick={0,2,4,8,16,32},tick label style={font=\footnotesize},grid=major,grid style={gray!25},
  table/col sep=comma,legend style={font=\scriptsize,at={(0.97,0.97)},anchor=north east}]
\addplot[omThm,thick,mark=*,mark size=1.4pt] table[x=gap,y=orders]{figdata/hft/29-a-venue-playbook/budget.csv};
\addplot[omDef,thick,dashed,mark=square*,mark size=1.3pt] table[x=gap,y=accepted]{figdata/hft/29-a-venue-playbook/budget.csv};
\legend{orders sent,orders accepted}
\end{axis}
\begin{axis}[width=11cm,height=5.4cm,axis y line*=right,axis x line=none,ymin=0,ymax=12,xmin=-1,xmax=33,
  ylabel={\small volume (thousand units)},tick label style={font=\footnotesize},table/col sep=comma]
\addplot[black,thick,dotted,mark=triangle*,mark size=1.6pt] table[x=gap,y=volume]{figdata/hft/29-a-venue-playbook/budget.csv};
\node[font=\scriptsize,anchor=south] at (axis cs:8,10.1) {volume (right)};
\end{axis}
\end{tikzpicture}

\omcaption{The crypto venue's order budget (one order every eight seconds, a burst of five) against the quoter's pacing: orders sent and accepted
(left, log scale) and volume traded (right), three seeds of 20 minutes. Up to an eight-second gap the venue accepts 151 to 154 orders whatever is
sent, the whole budget; volume peaks at eight seconds. Data: \texttt{hf\_venues.budget\_curve}.}\label{fig:hf:a-venue-playbook:budget}
\end{omfigure}

The budget binds in every run up to an eight-second gap: the venue accepts 151 to 154 orders in twenty minutes, 0.125 a second plus the burst,
whatever the quoter sends. What the pacing decides is which ones. A quoter that fires at every book update spends the budget in bursts and has no
orders left when the market moves; a paced one spends it evenly. Volume peaks at an eight-second gap, 10\,100 units, with 28 orders rejected in
bursts; the sixteen-second rule gives up a sixth of that volume to have no rejections at all. The cancels are free in this budget, as they are on the
venue in the dated box; the jump cancel never waits, which is why the profile's validation requires it.

An on-chain venue has no matching engine and no profile field for one: its fields are the block time, the pool's fee, the gas and the builder
auction that orders the transactions in a block (\cref{ch:hf:on-chain-trading}). The simulator does not run it; the profile records it so that the
desk's venue list is complete and so that a strategy that assumes an order book is refused there.

An event venue that delays new orders in play, and not cancels, protects resting orders: the arbitrageur's order waits five seconds, the quoter's
cancel does not, and only 5.6\% of the quoter's volume was picked off, against 42\% on the lit venue. But the delay applies to the quoter's own new
orders too, and 52\% of its volume there was taken, not made: an order priced for the book it saw lands five seconds later in a book that has
moved, crosses it, and trades as the aggressor at the old price. The edge was the worst of the eight venues after the batch, $-0.52$ ticks a unit. The profile's
parameter, a three-tick requote threshold, cut the messages by 70\% (799 to 238) for 6\% less volume, because a requote that lands five seconds
late is mostly wasted, and cut the taken share to 34\%; the edge did not improve ($-0.57$). Fixing it needs code that prices for where the book will
be when the order lands, or sends no new order while the price moves.

\section{Onboarding a new venue}

The batch venue is the other one no parameter fixes. It protects the quoter completely, no volume picked off, because the jump cancel lands in the
same interval as the arbitrageur's order and is processed before the batch clears. But a quoter that joins the touch in a uniform-price auction is
filled only when the batch clears through its price, which happens when the order flow in that batch runs against it: 2\,567 units at $-1.12$ ticks a
unit, every fill in the auction. The venue needs a quoter that prices against its fair value (\cref{ch:hf:fair-value}),
not the touch.

That result sets the order of onboarding. A desk writes the venue half of the profile from the rulebook, the fee schedule, the feed and
order-entry specifications and the technical notices; validates it with the home venue's quoting half, which lists what the home parameters break;
runs the strategy on the simulator configured from the profile, with the home and the proposed parameters, and reads the difference; then connects
to the venue's test environment, runs the conformance tests the venue requires, and starts with small size under the risk controls of
\cref{ch:hf:risk-controls}, as a staged rollout (One Quant Book 7, chapter 21).

\begin{center}
\small
\begin{tabular}{@{}lrrrrrr@{}}
\toprule
venue & volume & own & picked off & edge (ticks) & messages & own\\
\midrule
lit equity & 14\,567 & = & 42\% & $+0.03$ & 591 & =\\
inverted equity & 14\,567 & = & 42\% & $+0.03$ & 591 & =\\
bumped equity & 14\,567 & = & 42\% & $+0.03$ & 591 & =\\
batch & 2\,567 & = & 0\% & $-1.12$ & 330 & =\\
pro-rata futures & 10\,884 & 40\,633 & 47\% / 58\% & $-0.17$ / $-0.14$ & 474 & 419\\
FIFO futures & 14\,567 & = & 42\% & $+0.03$ & 591 & =\\
crypto & 7\,267 & 8\,333 & 56\% / 48\% & $-0.07$ / $-0.17$ & 2\,857 & 191\\
event & 8\,367 & 7\,900 & 6\% / 4\% & $-0.52$ / $-0.57$ & 799 & 238\\
\bottomrule
\end{tabular}
\end{center}

The table is the chapter's named result: the quoter's volume, the share picked off, its edge (ticks a unit against the efficient price five seconds
later) and its messages on each venue, with the home venue's parameters and, where they differ, the venue's own (an equals sign: nothing to
change). In this market joining the touch is roughly break-even on the continuous venues, as chapter 1 found; the venues differ in what takes the
edge away, and a profile's parameters recover only what parameters can.

The profile is kept alive after onboarding. Venues change fee schedules monthly, add order types, move data centres and change limits; each change is
a new version of the venue half, validated against the quoting half, and replayed through the simulator before it takes effect. The profile's date
is the date somebody last read the rulebook.

\section{Tutorial: one code, eight venues}\label{tut:hf:a-venue-playbook:tut}

\begin{tutorial}
\textbf{Goal.} Encode eleven \omterm{def:hf:a-venue-playbook:profile}{venue profiles}, validate each with the home venue's quoting parameters, and run one quoter on the eight simulated venues
with the home and the venue's own parameters. \textbf{End state:} the result table and the three figures.
\begin{tutsteps}
\item \textbf{The profiles}: \texttt{firm.venueprofile.standard()}, a \texttt{Registry} that validates each profile on entry and round-trips through
JSON.
\item \textbf{The rules}: \texttt{validate}; \texttt{hf\_venues.home\_errors} lists what the home parameters break (the crypto budget, the FX lot).
\omcode{code/firm/venueprofile/firm_venueprofile.py}{133}{153}{The rules that tie the quoting half to the venue half: lot, post-only, asymmetric delays and the order budget.}\label{lst:hf:a-venue-playbook:validate}
\item \textbf{The quoter}, written once: it joins the others' best prices, requotes past a threshold and no faster than its pacing, and cancels
its stale side on a jump signal.
\omcode{code/firm/venueprofile/firm_venueprofile.py}{303}{341}{The jump signal empties the stale side for a cooling period; the requote keeps one order a side at the target unless the target moved by the threshold.}\label{lst:hf:a-venue-playbook:quoter}
\item \textbf{The runs}: \texttt{hf\_venues.playbook} (three seeds of 20 minutes, home and own) and \texttt{hf\_venues.budget\_curve}.
\end{tutsteps}
\textbf{What to change next.} Give the batch venue a quoter that prices against a fair value; add a delay-aware quoter for the event venue that
sends no new order while the price moves; add the options venue with a mass-quote throttle.
\end{tutorial}

\section{Build: venue profiles}\label{bld:hf:a-venue-playbook:venueprofile}

\begin{build}
\bfield{Purpose} Keep one validated, versioned page per venue, and turn it into the simulator's configuration, the quoter's parameters and the
registry and fee-schedule rows.
\bfield{Interface} \texttt{Quoting}, \texttt{Profile} (\texttt{to\_dict}, \texttt{from\_dict}), \texttt{validate(p)}, \texttt{Registry} (\texttt{get},
\texttt{by\_type}, \texttt{simulated}, \texttt{dump}, \texttt{load}), \texttt{standard()}, \texttt{HOME}, \texttt{exchange\_config(p)},
\texttt{venue\_record(p)} (Book 1's \texttt{firm.venues}), \texttt{fee\_schedule(p)} (Book 1's \texttt{firm.feesched}), \texttt{jumps},
\texttt{Quoter}, \texttt{Sniper}, \texttt{run(p, quoting, seconds, seed)}.
\bfield{Rules} Prices in 1/10\,000 of the currency, as the simulator's; units as the market's; a profile that fails validation cannot enter a
registry; the budget applies to the quoter's session only (the background stands for many participants with their own).
\bfield{Acceptance tests} \texttt{code/firm/venueprofile/tests/}: every standard profile validates and the home parameters break the crypto budget
and the FX lot; a malformed profile lists all its faults; JSON round trip; the simulator's configuration for pro rata, batch, event and crypto; short
runs in which the bumped venue changes no race, the inverted venue changes the fee sign only, the event venue's delay protects resting orders, the
paced quoter has no rejects, and the larger pro-rata size trades more than twice the volume.
\bfield{Stretch} A profile diff that lists which rules a venue notice breaks; the options venue with \omterm{def:hf:automated-options-market-making:mass}{mass quotes}; versioned profiles with effective
dates.
\end{build}

\begin{omsources}
\item Nasdaq, TotalView-ITCH 5.0 specification.
\item US Securities and Exchange Commission, Release No.\ 34-78101 (IEX, 2016) and the order disapproving Cboe EDGA's liquidity provider protection
delay (February 2020).
\item CME Group, Globex Messaging Efficiency Program, rule filing with the CFTC, 8 November 2019.
\item Betfair Developer Program, ``Why do you have a delay on placing bets on a market that is in-play''.
\end{omsources}

\section{Exercises}

\begin{exercise}[$\star$]\label{exo:hf:a-venue-playbook:1}
The crypto budget is 0.125 orders a second with a burst of five. How many orders can the venue accept in twenty minutes, and how does that compare
with the runs?
\end{exercise}

\begin{exercise}[$\star$]\label{exo:hf:a-venue-playbook:2}
A requote sends at most one new order on each side. What minimum gap between requotes does the budget require, and what does the home venue's
quoting set?
\end{exercise}

\begin{exercise}[$\star$]\label{exo:hf:a-venue-playbook:3}
On the pro-rata book, an incoming order of 1\,000 contracts meets a level of 5\,000 resting contracts, of which the order that set the price holds
200 and the quoter 500. With a 40\% top-order share and the rest pro rata, how many does the quoter get?
\end{exercise}

\begin{exercise}[$\star\star$]\label{exo:hf:a-venue-playbook:4}
Why does a symmetric 350-microsecond delay change none of the quoter's fills in the runs, and what would change that?
\end{exercise}

\begin{exercise}[$\star\star$]\label{exo:hf:a-venue-playbook:5}
On the event venue only 6\% of the quoter's volume is picked off, yet its edge is $-0.52$ ticks. Explain.
\end{exercise}

\begin{exercise}[$\star\star$]\label{exo:hf:a-venue-playbook:6}
On the batch venue the quoter is never picked off. Why does it lose more than a tick a unit, and what would you change?
\end{exercise}

\begin{exercise}[$\star\star\star$]\label{exo:hf:a-venue-playbook:7}
\emph{Coding.} With \texttt{hf\_venues.budget\_curve}, which pacing maximises volume, how many orders does it have rejected, and why does it beat the
sixteen-second rule?
\end{exercise}

\begin{exercise}[$\star\star\star$]\label{exo:hf:a-venue-playbook:8}
\emph{Find the flaw.} ``The new venue's profile passed validation and the simulator shows our usual edge, so we can go live at full size on Monday.''
\end{exercise}

\section{Problem: One Code, Eight Venues}

\begin{problem}[{Weekend problem --- one code, eight venues}]\label{pb:hf:a-venue-playbook:1}
A market-making desk that quotes on a US equity exchange is asked to take its quoter to seven other venue types.

\medskip\noindent\textbf{Part I --- The venues.}
\begin{enumerate}
\item Define a \omterm{def:hf:a-venue-playbook:profile}{venue profile} and its two halves.
\item List the fields of the venue half and, for each, a chapter of this book where it mattered.
\item What do the dated box's venues say about matching, delays and budgets?
\item Why are an on-chain pool and an FX last-look stream in the registry although the simulator does not run them?
\end{enumerate}

\medskip\noindent\textbf{Part II --- The rules.}
\begin{enumerate}[resume]
\item Give the validation rules that tie the quoting half to the venue half.
\item Which rules do the home venue's parameters break, and on which venues?
\item Why must the jump cancel be free under the budget?
\item How does the quoter find its own orders, and what goes wrong if the feed shows an order before its acknowledgement arrives?
\end{enumerate}

\medskip\noindent\textbf{Part III --- Measurements.}
\begin{enumerate}[resume]
\item Describe the market of the runs: the flow, the jumps, the arbitrageur, the quoter's signal.
\item Give the volume, share picked off, edge and messages on each venue with the home parameters.
\item Which venues change with the venue's own parameters, and by how much?
\item Describe the budget curve and explain why the accepted orders are flat up to eight seconds.
\end{enumerate}

\medskip\noindent\textbf{Part IV --- The verdict.}
\begin{enumerate}[resume]
\item State the \emph{named result}: the quoter's capture and message count on each venue type with and without the venue-specific parameters.
\item Which venues need a parameter, which need nothing, and which need different code?
\item Why does the bumped venue not protect the quoter, and what kind of delay would?
\item What would you do on the event venue?
\item What would you do on the batch venue?
\item How would you onboard a ninth venue, step by step?
\item How do you keep a profile current, and who signs off a change?
\item In one sentence: what does a \omterm{def:hf:a-venue-playbook:profile}{venue profile} buy?
\end{enumerate}
\end{problem}

\section{Interview questions}

\begin{interviewq}[$\star$ \iqroles{developer}]\label{iq:hf:a-venue-playbook:1}
What fields would you put in a \omterm{def:hf:a-venue-playbook:profile}{venue profile}, and which three would you check before the first order?
\end{interviewq}

\begin{interviewq}[$\star\star$ \iqroles{trader}]\label{iq:hf:a-venue-playbook:2}
You move a quoting strategy from a price-time equities venue to a pro-rata futures book. What do you change first, and what is the risk of the
change?
\end{interviewq}

\begin{interviewq}[$\star\star$ \iqroles{developer}]\label{iq:hf:a-venue-playbook:3}
Your orders on a crypto venue are being rejected in bursts, and your quotes are missing when the market moves. Diagnose and fix.
\end{interviewq}

\begin{interviewq}[$\star\star$ \iqroles{researcher}]\label{iq:hf:a-venue-playbook:4}
A venue introduces a delay on incoming orders that spares cancels. How does it change a maker's economics, and whose?
\end{interviewq}

\begin{interviewq}[$\star\star$ \iqroles{researcher, trader}]\label{iq:hf:a-venue-playbook:5}
With an order budget of $b$ a second and a burst of $B$, and requotes of two orders each, what pacing uses the budget without rejects, and why might
a faster pacing trade more?
\end{interviewq}

\begin{interviewq}[$\star\star\star$ \iqroles{researcher}]\label{iq:hf:a-venue-playbook:6}
In a uniform-price batch auction a maker joins the best bid and offer. Show that its fills are adversely selected, and describe a quoting rule
that is not.
\end{interviewq}
